% 三元素底集 V={1,2,3} 的8个子集与指示向量 (x_1,x_2,x_3) 的对应：立方体顶点。
% 边连接只差一个元素（只差一个坐标）的两个子集；颜色表示改变的是哪个坐标。
\begin{tikzpicture}[line width=.8pt,font=\figurefont\fontsize{8}{10}\selectfont,
  v/.style={draw=FigureStroke,fill=FigureBlueFill,rounded corners=1.8pt,align=center,inner sep=2pt,minimum width=13mm,minimum height=8mm,line width=.8pt}]
  \def\ex{3.4} \def\dx{1.5} \def\dy{1.15} \def\ez{2.5}
  \foreach \a in {0,1}{\foreach \b in {0,1}{\foreach \c in {0,1}{
    \coordinate (p\a\b\c) at ({\a*\ex+\b*\dx},{\c*\ez+\b*\dy});
  }}}
  % 边：改变x1蓝、x2绿、x3黄
  \foreach \b in {0,1}{\foreach \c in {0,1}{\draw[FigureBlue,line width=1.1pt] (p0\b\c)--(p1\b\c);}}
  \foreach \a in {0,1}{\foreach \c in {0,1}{\draw[FigureGreen,line width=1.1pt] (p\a0\c)--(p\a1\c);}}
  \foreach \a in {0,1}{\foreach \b in {0,1}{\draw[FigureYellow,line width=1.1pt] (p\a\b0)--(p\a\b1);}}
  \node[v] at (p000) {$\varnothing$\\$(0,0,0)$};
  \node[v] at (p100) {$\{1\}$\\$(1,0,0)$};
  \node[v] at (p010) {$\{2\}$\\$(0,1,0)$};
  \node[v] at (p001) {$\{3\}$\\$(0,0,1)$};
  \node[v] at (p110) {$\{1,2\}$\\$(1,1,0)$};
  \node[v] at (p101) {$\{1,3\}$\\$(1,0,1)$};
  \node[v] at (p011) {$\{2,3\}$\\$(0,1,1)$};
  \node[v,fill=FigureYellowFill] at (p111) {$\{1,2,3\}$\\$(1,1,1)$};
  % 图例
  \begin{scope}[shift={(6.6,0.2)},font=\FigureNoteFont\fontsize{7.5}{9}\selectfont]
    \draw[FigureBlue,line width=1.1pt] (0,2)--(0.6,2) node[right,FigureInk]{加入或删去元素$1$};
    \draw[FigureGreen,line width=1.1pt] (0,1.5)--(0.6,1.5) node[right,FigureInk]{加入或删去元素$2$};
    \draw[FigureYellow,line width=1.1pt] (0,1)--(0.6,1) node[right,FigureInk]{加入或删去元素$3$};
  \end{scope}
\end{tikzpicture}
